\section{Quantitative defects of regular collision phases}
\label{sec:quantitative-phase-defects}
Absence of exact phases does not itself control approximate vectors.
We give a quantitative version on the space of bounded functions of
bounded variation. The phase is reconstructed from a complex vector;
no pointwise lower bound for its modulus is assumed. The logarithmic
bounds in this section are at a fixed radius and on a fixed compact
frequency set. A separate compactness argument gives a parameter-uniform
bound when the regularity budget is fixed.

For $\omega\in\Omega$ put
\begin{equation}\label{eq:collision-phase-defects}
 \begin{split}
 d_{R,\omega}(q)&=\|q\circ T_R-e^{i\phi_{R,\omega}}q\|_{L^1(\nu)},
                      \qquad (|q|=1),\\
 \mathcal E_{R,\omega}(f)&=\|f\circ T_R-e^{i\phi_{R,\omega}}f\|_{L^2(\nu)}.
 \end{split}
\end{equation}
The variation is always in the original collision coordinates. In
particular it is neither an anisotropic distribution norm nor a norm
of a many-return pushforward density.

\subsection{Finite-time measurable holonomy with a defect}
Use the product coordinates and the pair probabilities
\eqref{eq:explicit-conditional-pair-densities}. For a stable or unstable
oriented arc from $x$ to $y$ define
\[
 H_{s/u}(x,y)=
 \left|\frac{q(y)}{q(x)}-
            \exp\!\left(it\int_{\gamma_{s/u}(x,y)}\vartheta_R\right)\right|.
\]
These are bounded measurable functions of pairs. Null sets are taken
with respect to the indicated pair probabilities.

\begin{lemma}[A finite-time holonomy defect]
\label{lem:finite-holonomy-defect}
Fix a regular product set $\mathcal P$ at radius $R$. There are
$C_{\mathcal P}<\infty$ and $0<\theta_{\mathcal P}<1$ such that,
for $m\ge1$, $0<\epsilon<1/2$, and $q:M\to\mathbb S^1$ with
$\|q\|_{\mathrm{BV}}\le H$, $H\ge2$,
\begin{equation}\label{eq:finite-holonomy-defect}
 \int H_{s/u}\dd\mathfrak m_{s/u}
 \le C_{\mathcal P}\left[
 m d_{R,\omega}(q)+H\epsilon+
                  (\epsilon^{-1}+|t|)\theta_{\mathcal P}^{m}\right].
\end{equation}
The same expression controls both signs of time. The constants do not
depend on $u,s,t,H,m,\epsilon$.
\end{lemma}
\begin{proof}
Iterating the approximate phase equation along one orbit introduces
an error at most the sum of its $m$ one-step absolute defects. The
stable-pair version therefore has errors at its two endpoints. Both
pair marginals equal $m_{\mathcal P}=\nu(\mathcal P)^{-1}\nu|_{\mathcal P}$.
Invariance of $\nu$ bounds the integrated sum of endpoint defects by
$2m\nu(\mathcal P)^{-1}d_{R,\omega}(q)$.

Let $q_\epsilon=\mathsf S_\epsilon q$. Positivity of smoothing gives
$|q_\epsilon|\le1$, and
$\|q-q_\epsilon\|_1\le C H\epsilon$,
$\|q_\epsilon\|_{C^1(\alpha,p)}\le C\epsilon^{-1}$.
The two endpoint approximation errors are bounded at every iterate
by $2\nu(\mathcal P)^{-1}\|q-q_\epsilon\|_1$.
The actual stable arc lengths are at most
$C_{\mathcal P}\theta_{\mathcal P}^m$ after $m$ forward iterates,
by the contraction estimate in the proof of
Lemma~\ref{lem:regular-collision-product}. Thus the smoothed endpoint
difference costs $C_{\mathcal P}\epsilon^{-1}\theta_{\mathcal P}^m$.
The terminal action integral in \eqref{eq:action-telescope-arc} costs
at most $C_{\mathcal P}|t|\theta_{\mathcal P}^m$.
The lattice and constant phases cancel along the stable pair exactly.
These estimates prove \eqref{eq:finite-holonomy-defect} for stable arcs.

For an unstable pair telescope backwards. Its defect is the sum of the
one-step defects at $T_R^{-j}x,T_R^{-j}y$, $1\le j\le m$; the same
marginal domination applies by invertibility and invariance. Use
\eqref{eq:backward-action-arc} and backward contraction. The holonomy
sign is again positive, as in
\eqref{eq:measurable-unstable-holonomy}. This proves the other case.
No invariant pair measure and no independent orbit assumption is used.
\end{proof}

We make explicit how a measure-class assertion can be used
quantitatively without replacing it by an unproved bounded-density
assertion. Write $\mu=\mu^u\otimes\mu^s$ and
$\dd m_{\mathcal P}=\rho\dd\mu$ in the product coordinates.
Let $\lambda_s=\mu^u\otimes\mu^s\otimes\mu^s$ and
$\lambda_u=\mu^u\otimes\mu^u\otimes\mu^s$ be the reference triple
probabilities, and set
\[
 w_s=\frac{\dd\lambda_s}{\dd\mathfrak m_s},\qquad
 w_u=\frac{\dd\lambda_u}{\dd\mathfrak m_u}.
\]
These derivatives are positive and finite almost everywhere. For
$0\le F\le2$ and either pair type,
\begin{equation}\label{eq:pair-density-truncation}
 \int F\dd\lambda_{s/u}
 \le K\int F\dd\mathfrak m_{s/u}
            +2\lambda_{s/u}\{w_{s/u}>K\}.
\end{equation}
The last probability tends to zero as $K\to\infty$. Its value is part
of the constant below; no uniform pointwise bound on $\rho$ is needed.

\begin{lemma}[An almost constant phase on the product set]
\label{lem:almost-constant-product-phase}
For every $\eta>0$ there is a finite $C_{\mathcal P,\eta}$ such that
one can choose $c\in\mathbb S^1$ with
\begin{align}\label{eq:almost-constant-product-phase}
 \int_{\mathcal P}|q-c|\dd\nu
 \le\eta+C_{\mathcal P,\eta}\bigl[
 |t|+m d_{R,\omega}(q)+H\epsilon
              +(\epsilon^{-1}+|t|)\theta_{\mathcal P}^m\bigr].
\end{align}
This holds for the same $q,H,m,\epsilon$ as in the preceding lemma.
\end{lemma}
\begin{proof}
The action integrals on the original arcs are uniformly bounded on
$\mathcal P$. Hence the stable and unstable pair averages of
$|q(y)-q(x)|$ are bounded by the right side of
\eqref{eq:finite-holonomy-defect} plus $C_{\mathcal P}|t|$.

Here are the density choices needed to pass to two arbitrary product
points. Choose $J$ so that
$\int_{\{\rho>J\}}\rho\dd\mu$ is less than a prescribed small number.
Choose $K$ so that both tail probabilities in
\eqref{eq:pair-density-truncation} are less than that number divided
by $J$. Joining $(u,s)$ to $(u,s')$ and then to $(u',s')$, Fubini
bounds the $\mu\otimes\mu$ average of $|q(x)-q(y)|$ by the sum of
the two reference triple averages. Thus there is a value $c=q(y)$
for which the $\mu$ average of $|q-c|$ is at most that sum. Finally
\[
 \int|q-c|\dd m_{\mathcal P}
 \le J\int|q-c|\dd\mu+2\int_{\{\rho>J\}}\rho\dd\mu.
\]
Take the prescribed small numbers sufficiently small in terms of
$\eta$ and $\nu(\mathcal P)$, and apply
\eqref{eq:pair-density-truncation}. The remaining multiplier is a
finite function of $J,K,C_{\mathcal P},\nu(\mathcal P)$.
This proves \eqref{eq:almost-constant-product-phase}.
\end{proof}

\subsection{A quantitative roof obstruction}
\begin{proposition}[A logarithmic defect bound on a roof band]
\label{prop:logarithmic-roof-defect}
Fix $R$ and $0<t_0\le T<\infty$. There is $c_{R,t_0,T}>0$ such that
\begin{equation}\label{eq:logarithmic-roof-defect}
 d_{R,(u,s,t)}(q)\ge\frac{c_{R,t_0,T}}{1+\log H}
\end{equation}
whenever $t_0\le|t|\le T$, $|q|=1$, $H\ge2$, and
$\|q\|_{\mathrm{BV}}\le H$. The estimate is independent of $u,s$.
\end{proposition}
\begin{proof}
Restrict a regular product set so that every oriented four-corner
quadrilateral has action $\mathcal A=\oint\vartheta_R$ satisfying
$|\mathcal A|\le\pi/T$. By Stokes' formula and nonatomicity,
$|\mathcal A|>0$ for almost every quadruple under
$\lambda_4=(\mu^u)^2\otimes(\mu^s)^2$.
Choose $a_0>0$ such that
\[
 p_0:=\lambda_4\{|\mathcal A|\ge a_0\}>0.
\]
For this set and the entire roof band,
$|e^{it\mathcal A}-1|\ge2t_0a_0/\pi$.

Each of the four edge marginals of $\lambda_4$ is one of the reference
triple measures. Choose $K$ so that the sum of their probabilities
of $w_{s/u}>K$ is less than $p_0/2$. Intersect the set
$\{|\mathcal A|\ge a_0\}$ with all four complementary conditions;
call the resulting set $\mathcal Q$. It has measure at least $p_0/2$,
and each of its unnormalized edge marginals is bounded by
$K\mathfrak m_s$ or $K\mathfrak m_u$.

Multiply the four oriented approximate holonomies. The product of
$q$-ratios is exactly one, and the error in a product of unit complex
numbers is at most the sum of its four errors. Integration over
$\mathcal Q$ and Lemma~\ref{lem:finite-holonomy-defect} give
\begin{equation}\label{eq:quantitative-area-defect}
 \frac{p_0t_0a_0}{\pi}
 \le4K C_{\mathcal P}\bigl[
 m d_{R,\omega}(q)+H\epsilon
                  +(\epsilon^{-1}+T)\theta_{\mathcal P}^m\bigr].
\end{equation}
All constants have now been selected from the product set, its density
tails and its action distribution, independently of $q$ and $H$.
Choose $\epsilon=c_1/H$, with a fixed sufficiently small $c_1>0$.
Then choose $m=\lceil c_2+c_3\log H\rceil$ so that the final term in
brackets is also less than one quarter of the required lower bound.
The positive contraction exponent makes this possible. Equation
\eqref{eq:quantitative-area-defect} leaves $m d_{R,\omega}(q)\ge c_4>0$.
This proves \eqref{eq:logarithmic-roof-defect}.
\end{proof}

\subsection{All compact nonzero physical frequencies}
\begin{theorem}[A logarithmic phase defect on a compact band]
\label{thm:compact-log-phase-defect}
For a fixed $R\in I$ and a compact set
$\mathcal K\subset\Omega\setminus\{0\}$ there is $c_{R,\mathcal K}>0$
such that, for every $\omega\in\mathcal K$,
\begin{equation}\label{eq:compact-log-phase-defect}
 |q|=1,\quad \|q\|_{\mathrm{BV}}\le H,\quad H\ge2
 \quad\Longrightarrow\quad
 d_{R,\omega}(q)\ge\frac{c_{R,\mathcal K}}{1+\log H}.
\end{equation}
The set includes compact nonzero spatial frequencies with $s=t=0$.
\end{theorem}
\begin{proof}
Set
\[
 \rho(u,s)=|e^{iu_1}-1|+|e^{iu_2}-1|+|e^{is}-1|,
 \qquad \rho_*:=\min_{\mathcal K}(|t|+\rho(u,s))>0.
\]
Let $T=1+\max_{\mathcal K}|t|$. Fix a product set $\mathcal P$ and
the finite witnesses $(n_i,k_i),E_i,b_{ji}$ in
Lemma~\ref{lem:finite-return-witnesses}. Write $p_i=\nu(E_i)>0$ and
$d=d_{R,\omega}(q)$. For a constant $c\in\mathbb S^1$, iteration
of the approximate phase equation on $E_i$ gives
\begin{align}\label{eq:witness-phase-defect}
 p_i|e^{i(u\cdot k_i+s n_i)}-1|
 \le 2\int_{\mathcal P}|q-c|\dd\nu+n_i d
                     +p_i|t|n_i\|\tau_R\|_\infty.
\end{align}
Indeed $T_R^{n_i}E_i\subset\mathcal P$, its contribution to the
endpoint error is bounded by the integral over $\mathcal P$ by
invariance, and the $n_i$ one-step defects integrate to at most $n_i d$.
The roof sum is bounded by $n_i\|\tau_R\|_\infty$.

The integer identities $e_j=\sum_i b_{ji}g_i$ and
$|z^b-1|\le |b||z-1|$ for $|z|=1$, $b\in\Z$, show that finite
constants $A,B,C$ satisfy
\begin{equation}\label{eq:witness-resonance-control}
 \rho(u,s)\le A\int_{\mathcal P}|q-c|\dd\nu+B d+C|t|.
\end{equation}
For example one may take
$A=2\sum_{j,i}|b_{ji}|/p_i$,
$B=\sum_{j,i}|b_{ji}|n_i/p_i$ and
$C=\|\tau_R\|_\infty\sum_{j,i}|b_{ji}|n_i$.
These constants include the actual positive masses of the witnesses.

Choose $\eta>0$ with $A\eta<\rho_*/16$, and use
Lemma~\ref{lem:almost-constant-product-phase} with this $\eta$.
Choose $t_0\in(0,\rho_*/4)$ so small that all terms proportional
to $|t|$ in \eqref{eq:witness-resonance-control}, after substitution
of that lemma, are less than $\rho_*/16$ when $|t|\le t_0$.
Choose $\epsilon=c_5/H$, then $m=\lceil c_6+c_7\log H\rceil$,
so that the two remaining approximation terms contribute less than
$\rho_*/8$. For $|t|\le t_0$, the left side is at least
$3\rho_*/4$. The inequalities thus force $m d\ge c_8>0$.
For $|t|\ge t_0$ use Proposition~\ref{prop:logarithmic-roof-defect}
on $[t_0,T]$. Taking the smaller of the two constants proves the
assertion. All selections precede the choice of $q,H$.
\end{proof}

\subsection{Reconstructing a phase without a lower-modulus assumption}
\begin{lemma}[A bounded-variation circle truncation]
\label{lem:bv-circle-truncation}
For every complex $f\in\mathrm{BV}(M)$ there is a measurable
$q:M\to\mathbb S^1$ such that
\begin{equation}\label{eq:bv-circle-truncation}
 \|q\|_{\mathrm{BV}}\le C(1+\|f\|_{\mathrm{BV}}),\qquad
 \|q-f\|_2\le3\||f|-1\|_2
\end{equation}
whenever the latter norm is finite.
\end{lemma}
\begin{proof}
Put $a=|f|$. The BV chain inequality gives $\TV(a)\le C\TV(f)$.
The coarea formula gives a level $\tau\in[1/4,1/2]$ for which
$E=\{a>\tau\}$ has perimeter at most $4\TV(a)$.
These chain, product and coarea statements are for first distributional
derivatives in the collision coordinates; see \cite[Chapter 3]{AFP}.
Let $N_\tau(z)=z/\max\{|z|,\tau\}$ and define
$q=N_\tau(f)\one_E+\one_{M\setminus E}$.
The map $N_\tau$ is uniformly Lipschitz for the selected levels and
has modulus at most one. The BV product inequality, including the
jump on the finite-perimeter boundary of $E$, proves the first bound.
On $E$, $|q-f|=|1-a|$. On its complement $a\le1/2$ and
$|q-f|\le1+a\le3(1-a)$. This proves the second bound.
In particular no value of $f/|f|$ is used at a zero of $f$.
\end{proof}

\begin{theorem}[A logarithmic obstruction for complex approximate vectors]
\label{thm:compact-vector-defect}
Fix $R$ and a compact $\mathcal K\subset\Omega\setminus\{0\}$.
There is $c'_{R,\mathcal K}>0$ such that
\begin{equation}\label{eq:compact-vector-defect}
 \|f\|_2=1,\quad \|f\|_\infty+\|f\|_{\mathrm{BV}}\le H,\quad H\ge2
 \quad\Longrightarrow\quad
 \inf_{\omega\in\mathcal K}\mathcal E_{R,\omega}(f)
             \ge\frac{c'_{R,\mathcal K}}{(1+\log H)^2}.
\end{equation}
The bound applies to complex $f$, with no restriction on its zero set.
\end{theorem}
\begin{proof}
Write $\varepsilon=\mathcal E_{R,\omega}(f)$, $a=|f|$,
$\bar a=\int a\dd\nu$ and $x=\|a-\bar a\|_2$.
The reverse triangle inequality gives
$\|a\circ T_R-a\|_2\le\varepsilon$, hence
$\|a\circ T_R^m-a\|_2\le m\varepsilon$.
The BV correlation bound of Proposition~\ref{prop:bv-correlation},
rescaled by homogeneity, gives constants $C_0,\gamma>0$ with
\begin{equation}\label{eq:approximate-modulus-variance}
 x^2\le C_0 H^2e^{-\gamma m}+x m\varepsilon.
\end{equation}
Indeed add and subtract the lag-$m$ covariance of the centered $a$ and
use Cauchy--Schwarz on the remaining term. The constants in this
collision estimate can be taken uniform in $R$.

Given $0<\eta<1$, choose
$m\le C[1+\log H+|\log\eta|]$ with
$C_0H^2e^{-\gamma m}\le\eta^2$.
Equation~\eqref{eq:approximate-modulus-variance} implies
$x\le\eta+m\varepsilon$. Since $\|a\|_2=1$ and $\bar a\ge0$,
\[
 \|a-1\|_2^2=2(1-\bar a)
       \le2(1-\bar a^2)=2x^2.
\]
Lemma~\ref{lem:bv-circle-truncation} constructs $q$ with
$\|q\|_{\mathrm{BV}}\le C(1+H)$ and
$\|q-f\|_2\le3\sqrt2(\eta+m\varepsilon)$. Invariance and the
unit modulus of the multiplier yield
\begin{equation}\label{eq:phase-reconstruction-defect}
 d_{R,\omega}(q)
 \le\mathcal E_{R,\omega}(q)
 \le\varepsilon+6\sqrt2(\eta+m\varepsilon).
\end{equation}
This is an explicit reconstruction and defect comparison.

Let $L=1+\log H$. Theorem~\ref{thm:compact-log-phase-defect}, with
the bound $C(1+H)$ for $q$, makes the left side at least $c_9/L$.
Choose $\eta=c_{10}/L$ so that $6\sqrt2\eta<c_9/(2L)$.
The selected $m$ is at most $C_{R,\mathcal K}L$, since
$\log L\le L$. Equation~\eqref{eq:phase-reconstruction-defect}
then gives $\varepsilon\ge c'_{R,\mathcal K}/L^2$.
No lower-modulus estimate for $f$ was assumed; it was replaced by
\eqref{eq:approximate-modulus-variance} and the controlled truncation.
\end{proof}

\begin{corollary}[Polynomial regularity budgets on a fixed band]
\label{cor:polynomial-bv-vector-budget}
For fixed $R,\mathcal K,H_0,\zeta$ with $H_0\ge2$, $\zeta\ge0$,
normalized vectors with
$\|f_N\|_\infty+\|f_N\|_{\mathrm{BV}}\le H_0N^\zeta$
satisfy $\inf_{\omega\in\mathcal K}\mathcal E_{R,\omega}(f_N)
\ge c/(1+\log(2+N))^2$. In particular no such vectors have defect
$O(N^{-a})$ for any fixed $a>0$.
\end{corollary}
\begin{proof}
Substitute the regularity budget into
\eqref{eq:compact-vector-defect}. A positive power of $N$ dominates
$(1+\log(2+N))^2$.
\end{proof}

\subsection{Uniform compact separation across the physical family}
\begin{theorem}[Parameter-uniform separation at bounded regularity]
\label{thm:uniform-compact-bv-separation}
For every $H<\infty$ and compact
$\mathcal K\subset\Omega\setminus\{0\}$ there is
$c_{H,\mathcal K}>0$ such that
\begin{equation}\label{eq:uniform-compact-bv-separation}
 \inf_{\substack{R\in I,\ \omega\in\mathcal K,\ \|f\|_2=1\\
                \|f\|_\infty+\|f\|_{\mathrm{BV}}\le H}}
          \mathcal E_{R,\omega}(f)\ge c_{H,\mathcal K}.
\end{equation}
An empty admissible class causes no difficulty. The theorem asserts
no uniform logarithmic dependence of this constant on $H$.
\end{theorem}
\begin{proof}
Suppose a sequence made the infimum zero. Pass to subsequences with
$R_j\to R$, $\omega_j\to\omega\in\mathcal K$ and $f_j\to f$ in
$L^1(\nu)$. The last subsequence exists by BV compactness. In the present
coordinates it also follows directly from
Lemma~\ref{lem:bv-smoothing}: the approximation is uniform in $L^1$,
and at every fixed smoothing scale the bounded smoothed functions
form a precompact family on the reflected compact cylinder.
The supremum bound upgrades this convergence to $L^2$, preserves
$\|f\|_2=1$, and gives $|f|\le H$ almost everywhere.

Outside the finite one-collision singularity set at $R$,
$T_{R_j}x\to T_Rx$, $\kappa_{R_j}(x)\to\kappa_R(x)$ and
$\tau_{R_j}(x)\to\tau_R(x)$, by persistence of the regular physical
root. Thus $e^{i\phi_{R_j,\omega_j}}\to e^{i\phi_{R,\omega}}$
almost everywhere and in every finite $L^p$ norm. For a smooth $g$,
invariance of the common probability gives
\begin{align*}
 \|f_j\circ T_{R_j}-f\circ T_R\|_2
 \le\|f_j-f\|_2+2\|f-g\|_2
                      +\|g\circ T_{R_j}-g\circ T_R\|_2.
\end{align*}
Take $j\to\infty$, then $g\to f$ in $L^2$. Dominated convergence
controls the last term. The multiplier term also converges in $L^2$.
The limiting defect is therefore zero.

The identity for $f$ makes $|f|$ invariant. Ergodicity and $\|f\|_2=1$
give $|f|=1$ almost everywhere. Theorem~\ref{thm:complete-physical-phase}
now forces $\omega=0$, contradicting $\mathcal K\subset\Omega\setminus\{0\}$.
This proves the positive infimum, without claiming continuity of
product sets or their density-tail constants in $R$.
\end{proof}

\paragraph{Relation to the complementary integral.}
Theorems~\ref{thm:compact-vector-defect} and
\ref{thm:uniform-compact-bv-separation} are estimates on actual
collision functions, with their stated norms. They cannot be applied
to an anisotropic distribution until a map to such functions, its
normalization, regularity and defect comparison have been proved.
They also do not assert that these BV classes are invariant under
the twisted collision operator. Absence of exact phases alone has not
been substituted for either assertion.

The logarithmic estimate allows a growing regularity budget on a fixed
compact nonzero band at each fixed $R$. The parameter-uniform theorem
instead fixes that budget. No estimate of $c_{R,\mathcal K}$ is supplied
when the roof band grows or $\mathcal K$ approaches the origin.
Consequently the physical annulus beginning at
$2n^{-99/200}$, equivalently the rescaled boundary $2n^{1/200}$,
is not declared controlled by these theorems. The full complementary
integral, the complete critical and singular raw residual sum and the
weighted exact-event comparison retain their distinct roles in
Theorem~\ref{thm:LLT}. The present estimates give a quantitative
phase and function-vector step toward that same inversion problem.
